\documentclass[11pt]{article}

\usepackage[margin=1in]{geometry}
\usepackage{amsmath,amssymb,amsthm,mathtools}
\usepackage{enumitem}
\usepackage{booktabs}
\usepackage{tabularx}
\usepackage{microtype}
\usepackage[hidelinks]{hyperref}
\usepackage[nameinlink,noabbrev]{cleveref}
\usepackage{xurl}

\newtheorem{definition}{Definition}
\newtheorem{remark}{Remark}

\title{Auditor Replay for Anonymous-DHT Receipt Line Items\\
\large A Mechanical Checklist for Digest-Bound Claims (Series Synthesis)}
\author{}
\date{Unified archive note v0.17 (2026-03-20; archive v453)}

\begin{document}
\maketitle

\begin{abstract}
This archive publishes privacy claims as compact receipts: an envelope plus line items with digest-bound replay hooks.
Many papers implicitly assume an auditor can ``replay the accountant'' from a receipt and an artifact bundle, but they do not spell out the mechanical steps.
This note is the canonical replay checklist.
It standardizes: (i) what objects an auditor needs, (ii) what must be checked before doing any privacy math, and (iii) what a successful replay verdict should record.
It is intentionally not a privacy proof; it is a verification workflow contract so other papers can stay short.
\end{abstract}

\paragraph{Scope.}
This note standardizes \emph{mechanical replay}: digest checks, envelope matching, log anchoring, and plan-spec replay.
It does \emph{not} define endpoint metrics, witness tiers, or composition theorems; those are imported.

\paragraph{Imports.}
Receipt schema (envelope + line items) is owned by Synthesis~12.\cite{series:receiptitems}
Plan/state drift and digest binding rules are owned by Synthesis~18.\cite{series:planstateequiv}
Horizon-$n$ accounting rules (retries/state) are owned by Synthesis~19.\cite{series:accountingrulebook}
Conditional per-step certificate shape for Regime~II is owned by Synthesis~21.\cite{series:condcert}
ABOM/OINL packaging is owned by Anonymity~C.\cite{series:abomoinl}
Transparency primitives (append-only logs, inclusion/consistency) are owned by Synthesis~9.\cite{series:transparencyprimitives}
Deployed-surface ownership (what belongs to state, to a service contact surface, or to the ENF/projection) is owned by Synthesis~30.\cite{series:deployedsurfaces}
A concrete end-to-end instantiation appears in the worked example (Synthesis~17).\cite{series:workedexample}

\paragraph{Exports.}
A checklist (\cref{sec:checklist}) and a small \emph{replay verdict} record format (\cref{sec:verdict}) that auditors and tooling can emit.

\paragraph{Later-terse-paper citation posture.}
Later terse papers should cite this note only when the live issue is the mechanical replay checklist itself, the replay-verdict field set, or the replay-derived successor continuity capsule. If the live issue is instead receipt tuple semantics, drift taxonomy / action-class classification, or deployed-surface assembly, they should stop at Synthesis~12 / Synthesis~18 / Synthesis~30 respectively. If the question has already widened to release-stage ownership---obligation profiles, package closure, public-status wrappers, or stable maintenance stages---they should widen to Synthesis~32 rather than treating replay as a release-spine note. Once a checked-answer handoff already exists and the only remaining question is whether a terse paper may stop, widen to the adjacent source-only branch, or reopen a named owner, they should cite Synthesis~74--75 instead of reopening replay by habit.\cite{series:receiptitems,series:planstateequiv,series:deployedsurfaces,series:releasespine,series:pairedterminalcitationdiscipline,series:pairedterminalcutoverrules}

\paragraph{A minimal public replay card.}
Later terse papers often need to say only that one declared tuple really was anchored, resolved, and rerun. The smallest honest public answer is therefore one replay card naming the exact audited receipt digest, the exact line item, the exact replay hook, the resolved state/contact surface ids, the checkpoint, and the recomputed scalar/verdict. That lets later notes stop at the replay result without republishing the whole checklist.

\begin{table}[t]
\centering\small
\begin{tabularx}{\linewidth}{@{}lX@{}}
\toprule
Slot & Minimal public content \\ 
\midrule
\texttt{receipt\_digest} / line-item key & Which exact receipt object and which exact tuple entry were replayed. \\ 
\texttt{plan\_id} / \texttt{plan\_spec\_id} & The stable replay recipe name and the digest-bound concrete plan specification. \\ 
Resolved surface ids & The exact \texttt{state\_decl\_id}, any \texttt{contact\_surface\_id}, and any adjunct used to assemble the effective surface. \\ 
\texttt{log\_checkpoint} & The inclusion/consistency checkpoint under which the replayed receipt was anchored. \\ 
Recomputed scalar(s) & The replayed privacy quantity and the verdict (pass/fail, plus any regime note if horizon accounting was used). \\ 
Optional continuity fields & Action-class / continuity / obligation-profile fields only when the replay is answering a successor-maintenance question. \\ 
\bottomrule
\end{tabularx}
\caption{A minimal public replay card. This is the smallest public answer later terse papers can quote directly when the live issue is \emph{did the declared tuple actually replay under the declared surface and checkpoint?} rather than the wider release-status or citation-spine story.}
\label{tab:minreplaycard}
\end{table}

\paragraph{Tiny fail-fast replay drill.}
Suppose a later note only needs to report: ``receipt digest $h(R)$ was anchored at checkpoint $c$, line item $(\texttt{claim\_id},\texttt{exposure\_nf\_id})$ replayed under $\texttt{plan\_spec\_id}=p$, and the recomputed bound stayed at $0.42$ bits/lookup.'' The honest short answer is Table~\ref{tab:minreplaycard} plus one sentence recording the resolved \texttt{state\_decl\_id} / \texttt{contact\_surface\_id} pair and the pass verdict. If the digest mismatches, if the resolved surface ids changed, or if the recomputed scalar must be interpreted through successor action classes and package duties, the terse paper has already widened beyond replay and should reopen Synthesis~18 or Synthesis~32 rather than pretending the replay card alone settles the release question.\cite{series:planstateequiv,series:deployedsurfaces,series:releasespine}

\section{Objects an auditor needs (minimal set)}
An auditor should be able to do a replay using only public objects plus an on-request artifact bundle.
The minimal set is:
\begin{itemize}[leftmargin=*]
\item \textbf{Receipt} $R$ (envelope + line items).\cite{series:receiptitems}
\item \textbf{ABOM} $A$ (manifest committing to claim semantics and artifact digests).\cite{series:abomoinl}
\item \textbf{UVI/log anchor} $U$ (receipt digest + inclusion/consistency proofs), or an equivalent transparency checkpoint chain.\cite{series:transparencyprimitives}
\item \textbf{Registries} used to bind human labels to digest ids (ENF/TW/State registries, when the receipt uses digest ids).
\item \textbf{On-request artifact bundle} referenced by each line item's \texttt{artifact\_bundle} (the omitted proofs, logs, scripts, certificates).
\end{itemize}

\begin{remark}[Why ABOM is an input, not optional]
In this archive, the receipt is intentionally small.
Without ABOM, an auditor cannot tell whether the artifact bundle presented on request is \emph{the} bundle the publisher committed to when publishing the receipt.
ABOM provides the digest commitments that make ``send me the artifacts'' checkable.\cite{series:abomoinl}
\end{remark}

\section{Mechanical replay checklist}\label{sec:checklist}
The checklist below is intentionally ordered: fail fast on anchoring and digest mismatch before doing any privacy-domain computation.

\subsection{Step 0: log anchoring and signature hygiene}
\begin{enumerate}[leftmargin=*]
\item \textbf{Compute receipt digest.} Compute $h(R)$ under the canonical digest rule the series uses.
\item \textbf{Verify inclusion.} Verify the Merkle inclusion proof that commits $h(R)$ under the stated checkpoint.\cite{series:transparencyprimitives}
\item \textbf{Verify consistency context.} Verify the checkpoint chain or consistency proof from the auditor's last checkpoint.\cite{series:transparencyprimitives}
\item \textbf{Verify signatures.} Verify the receipt/ABOM signatures under the declared keys (release key, notarization key, or both), as specified by the publication layer.\cite{series:abomoinl}
\end{enumerate}

\subsection{Step 1: envelope matching (drift guard)}
Receipts are compared across releases by their envelope.
Before replaying any line item, check that the published objects agree on the envelope IDs:
\begin{enumerate}[leftmargin=*]
\item \textbf{Receipt envelope.} Read \texttt{claim\_id}, \texttt{tw\_id}, and \texttt{state\_decl\_id} from the receipt envelope.\cite{series:receiptitems}
\item \textbf{ABOM claim fields.} Confirm ABOM/OINL pins the same claim family and threat-window semantics.
\item \textbf{UVI pointer.} Confirm the user-facing pointer (if present) names the same receipt digest and (at least) the same \texttt{claim\_id} and \texttt{exposure\_nf\_id}.\cite{series:transparencyprimitives}
\end{enumerate}
If any envelope ID differs, the auditor should treat the objects as distinct claims and avoid ``diffing'' budgets as if they were comparable.\cite{series:planstateequiv}

\subsection{Step 2: line-item integrity and replay-hook checks}
For each line item $\ell$ in $R$:
\begin{enumerate}[leftmargin=*]
\item \textbf{Projection binding.} Confirm \texttt{exposure\_nf\_id} is a digest id (or label resolvable to one) and matches the extractor/schema the mechanism paper claims to budget. In particular, reader-visible transcript semantics such as status-code interpretation, empty-result equivalence, chunking, or streaming shape belong here rather than in ad hoc prose.\cite{series:receiptitems,series:deployedsurfaces}
\item \textbf{Threat-window binding.} Confirm \texttt{tw\_id} matches the receipt envelope or is explicitly overridden (rare; should be justified).\cite{series:receiptitems}
\item \textbf{Replay hook is digest-bound.} Confirm \texttt{replay\_hook} includes \texttt{plan\_id} and \texttt{plan\_spec\_id}, and that \texttt{plan\_spec\_id}=\texttt{sha256(canon(plan\_spec))} for the on-request plan spec.\cite{series:planstateequiv}
\item \textbf{Resolve the effective surface.} Before replaying, assemble the actual service-facing interface from the declared \texttt{state\_decl\_id} (router-set precedence, AutoConf expansion, allow/deny rules), any line-item \texttt{contact\_surface\_id}, and the ENF-owned response semantics. The worked example now ships a tiny adjunct for this step so the rule is checkable without turning it into a new schema.\cite{series:deployedsurfaces,series:workedexample}
\item \textbf{Stop at replay unless the task widens.} If the question is still only ``can the declared computation be rerun?'' the replay anchor is the first sufficient owner. The maintained worked bundle now ships a tiny question-routing catalog that binds each shipped public answer to one public anchor, one stop-profile key, and one review-menu widening path, so auditors can distinguish replay from later release-status or referee-handoff questions instead of sliding rightward through the archive by habit.\cite{series:workedexample}
\item \textbf{State surface scope.} If the line item depends on state (explicitly or by regime), confirm the relevant \texttt{state\_decl\_id} is present (either at the envelope or inside the line item) and matches the declared state scope.\cite{series:accountingrulebook}
\item \textbf{Action-class lookup.} If the replay is motivated by a successor release, consult the release-action matrix or equivalent drift policy adjunct (when published) and record which field-family rule triggered the replay request. This keeps the audit verdict comparable to future releases without making the replay note itself own change-control prose.\cite{series:planstateequiv,series:workedexample}
\end{enumerate}

\subsection{Step 3: artifact-bundle commitment check}
For each referenced \texttt{artifact\_bundle}:
\begin{enumerate}[leftmargin=*]
\item \textbf{Resolve bundle.} Use the public bundle map (if provided) to list required artifacts for replay.
\item \textbf{Verify digests.} Confirm each delivered artifact matches the ABOM-committed digest list.
\item \textbf{Verify provenance.} If the plan spec requires notarization or transparency anchoring (e.g., a plan log inclusion proof), verify those proofs before using the artifact as evidence.\cite{series:transparencyprimitives}
\end{enumerate}

\subsection{Step 4: replay the accountant}
Finally, run the replay plan described by \texttt{plan\_spec}:
\begin{enumerate}[leftmargin=*]
\item \textbf{Recompute the published quantity.} Using the plan spec, recompute the line item's target scalar (e.g., per-lookup MaxL bound after attenuation, selection-tax term, or a TV dual-certificate bound).
\item \textbf{Check interpretation.} Confirm the receipt's \texttt{budget\_type} matches the endpoint metric meaning (imported from the endpoint bridge) and that the unit is consistent with \texttt{usage} (per lookup vs horizon).\cite{series:accountingrulebook}
\item \textbf{Record regime.} If a horizon claim is made (e.g., ``$Q\times$'' or a published total), record which regime of Synthesis~19 was used and which independence/state assumptions were relied on.\cite{series:accountingrulebook}
\end{enumerate}

\section{Replay verdict record (what to output)}\label{sec:verdict}
To keep audits diffable across time, an auditor should emit a small verdict record.
A minimal record is:
\begin{table}[t]
\centering\small
\begin{tabularx}{\linewidth}{@{}lX@{}}
\toprule
Field & Meaning \\
\midrule
\texttt{receipt\_digest} & Digest of the exact receipt object audited \\
\texttt{envelope} & \texttt{claim\_id}, \texttt{tw\_id}, \texttt{state\_decl\_id} \\
\texttt{line\_item\_verdicts} & For each line item: \texttt{exposure\_nf\_id}, \texttt{plan\_id}, \texttt{plan\_spec\_id}, recomputed scalar(s), and pass/fail \\
\texttt{resolved\_surface\_ids} & The \texttt{state\_decl\_id}, any \texttt{contact\_surface\_id}, and any adjunct used to assemble the effective replay surface \\
\texttt{release\_action\_class} & Optional: refresh/replay/recertify classification and the matched rule id when the replay responds to release drift rather than a fresh standalone audit \\
\texttt{continuity\_decision} & Optional: whether the audited successor remains on the same certified interface, and if so under what replay scope (for example, same-interface replayed-surface continuity) \\
\texttt{publication\_obligation\_profile} & Optional: the obligation-profile id and the required output roles when the replay is part of successor-release maintenance rather than a standalone audit \\
\texttt{artifact\_bundle\_digests} & Digests of the artifact bundles actually replayed \\
\texttt{log\_checkpoint} & The checkpoint used for inclusion/consistency \\
\texttt{notes} & Any regime/assumption notes (e.g., independence vs stateful accounting) \\
\bottomrule
\end{tabularx}
\caption{A minimal replay verdict is itself a small diffable object.}
\end{table}

\paragraph{Successor-comparison capsules.}
If the replay is performed because a specific successor release drifted, the auditor may publish a second tiny adjunct derived from the replay verdict: a \emph{successor continuity verdict}. That capsule simply records the base/successor pair, the matched release-action rule, the replay scope, the resulting continuity decision, and (optionally) the publication-duty roles imported from an obligation profile. It does not replace the replay verdict; it is the smallest base-to-successor conclusion object that later papers or release notes can cite without republishing the whole replay checklist. The worked example now ships one concrete instance.\cite{series:workedexample,series:planstateequiv}

\paragraph{Package-closure handoff.}
The replay note should stop one step short of owning release-package theory. Its role is only to emit enough structured information that a later maintenance adjunct can check package completeness. Concretely: if the replay was triggered by successor drift, the replay verdict and successor continuity verdict may carry the matched action class and required publication roles; a separate \emph{publication-closure verdict} can then bind those roles to the concrete shipped artifacts. This keeps Synthesis~20 about replay and evidence, while Synthesis~18 owns the classification/closure vocabulary.\cite{series:workedexample,series:planstateequiv}

\paragraph{Public-status handoff.}
The replay note should also stop short of owning the final release-note wrapper. If maintainers want a single public answer, a later adjunct such as a \emph{successor-lineage notice} may quote the matched action class, continuity decision, closure status, and replay-scope summary, then point readers back to the compare report and verdict objects. That wrapper is useful precisely because it is derived from replay and closure rather than replacing them. The replay note only needs to emit enough structured fields that such a wrapper can be checked.\cite{series:workedexample,series:planstateequiv}

\paragraph{Derivation-graph handoff.}
If maintainers also want one citeable statement of dependency order, a later adjunct such as a \emph{successor-derivation graph} may name the concrete compare report, replay verdict, successor continuity verdict, publication-closure verdict, and lineage notice, then connect them with a tiny DAG. That still keeps Synthesis~20 honest: replay exports evidence and scope, but a downstream maintenance note owns the graph that says how those pieces assemble into the final successor status.\cite{series:workedexample,series:planstateequiv}


\paragraph{Release-spine handoff.}
If maintainers also want one citeable statement of the \emph{stable maintenance stages} rather than the concrete DAG, a later adjunct such as a \emph{release spine} may name the canonical order diff $\to$ classify $\to$ obligate $\to$ close $\to$ summarize $\to$ support $\to$ cite $\to$ quote $\to$ clause, together with the smallest artifact answering each stage question for one base$\to$successor comparison. Replay still should not own that stage map. Its role is only to emit enough checked evidence and replay scope that the later spine can stay anchored to the narrower maintenance objects instead of drifting into release-note prose.\cite{series:workedexample,series:planstateequiv}

\paragraph{Downstream normal-form handoff.}
Replay should also stop short of owning sentence choreography.
Once the public-status wrapper already exists, the downstream suffix
\emph{summarize $\to$ support $\to$ cite $\to$ quote $\to$ clause}
belongs to Synthesis~31: claim-support maps, citation maps, quote maps, clause packs, and any tiny normal-form wrapper built on them.
In the current worked cut, that wrapper also carries typed leftward resolution traces, so replay remains the source of checked evidence while Synthesis~31 owns how an emitted sentence resolves back to it.\cite{series:downstreamnormalforms,series:workedexample,series:planstateequiv}

\begin{remark}[Why the verdict is part of the publication story]
A replay verdict is a way to publish \emph{that the replay was performed} without publishing the full private artifact archive.
It is compatible with the source-only posture: the verdict can be public and small; the artifacts remain available on request.
\end{remark}

\section{Worked example pointer}
The worked-example bundle (Synthesis~17) provides concrete instances of every checklist step: envelope IDs, registries, replay plans, support-bundle maps, a release-action matrix, a publication-obligation profile, a publication-closure verdict, a successor-lineage notice, a successor-derivation graph, a release spine, a successor-claim support map, a successor-citation map, a successor-quote map, a successor-clause pack, and a deterministic validation report.\cite{series:workedexample}
The purpose of this note is to make that workflow citable without requiring every paper to point directly at the worked example.

\begin{thebibliography}{99}\small

\bibitem{series:receiptitems}
Anonymous.
\newblock Receipt Line Items for Anonymity Claims: A Minimal Tuple Schema for Published Budgets.
\newblock Series synthesis note (Synthesis~12), 2026. (This archive.)

\bibitem{series:planstateequiv}
Anonymous.
\newblock Plan and State Equivalence: Digest-Bound Replay Hooks and Drift Taxonomy.
\newblock Series synthesis note (Synthesis~18), 2026. (This archive.)

\bibitem{series:accountingrulebook}
Anonymous.
\newblock Receipt Accounting and Horizon Composition for Anonymous-DHT Budgets.
\newblock Series synthesis note (Synthesis~19), 2026. (This archive.)

\bibitem{series:abomoinl}
Anonymous.
\newblock ABOM/OINL: Anonymity Bills of Materials and Odds-Inflation Nutrition Labels.
\newblock Anonymity series Paper~C, 2026. (This archive.)

\bibitem{series:transparencyprimitives}
Anonymous.
\newblock Transparency Primitives for Receipts and Certificates.
\newblock Series synthesis note (Synthesis~9), 2026. (This archive.)

\bibitem{series:condcert}
Anonymous.
\newblock Conditional Per-Step Budget Certificates: A Receipt-Grade Interface for Adaptive Horizon Claims.
\newblock Synthesis~21, The-One-True-Archive (2026).

\bibitem{series:workedexample}
Anonymous.
\newblock Worked Example: Receipt Interlock for an Anonymous-DHT Reader-Privacy Claim.
\newblock Series synthesis note (Synthesis~17), 2026. (This archive.)

\bibitem{series:deployedsurfaces}
Anonymous.
\newblock Deployed Observation Surfaces for Anonymous DHTs: Control-plane knobs in IPFS/libp2p delegated routing and OHTTP evolution.
\newblock Series synthesis note (Synthesis~30), 2026. (This archive.)

\bibitem{series:downstreamnormalforms}
Anonymous.
\newblock Downstream Normal Forms for Anonymous-DHT Release Notes.
\newblock Series synthesis note (Synthesis~31), 2026. (This archive.)

\bibitem{series:releasespine}
Anonymous.
\newblock Release Spine Contracts for Anonymous-DHT Successor Maintenance.
\newblock Series synthesis note (Synthesis~32), 2026. (This archive.)

\bibitem{series:pairedterminalcitationdiscipline}
Anonymous.
\newblock Paired Terminal Citation Discipline for Review Spines.
\newblock Series synthesis note (Synthesis~74), 2026. (This archive.)

\bibitem{series:pairedterminalcutoverrules}
Anonymous.
\newblock Paired Terminal Cutover Rules for Review Spines.
\newblock Series synthesis note (Synthesis~75), 2026. (This archive.)

\end{thebibliography}
\end{document}
